\clearpage
\section{Entwicklungs- und Fehlerchronologie}

Dieser Anhang bewahrt wesentliche negative Ergebnisse und Korrekturen. Im
Haupttext werden nur die wissenschaftlich entscheidenden Datenverträge und
Schutzmaßnahmen erläutert.

\begin{table}[H]
\centering
\scriptsize
\renewcommand{\arraystretch}{1.05}
\setlength{\tabcolsep}{3pt}
\begin{tabularx}{\textwidth}{p{2.6cm}XXX}
\toprule
Stufe & Beobachtung & Ursache und Korrektur & Erkenntnis\\
\midrule
SAM~3.1 & Speicherfehler oder formal leere Masken & Laufzeit und Parser wurden
an das verwendete Modell angepasst; Ausgaben werden unmittelbar persistiert
und über Stichproben sowie Coverage geprüft. & Ein erfolgreicher Prozess ist
kein ausreichender Maskennachweis.\\
COLMAP & Mehrere Sparse-Teilmodelle und ungeeignete Initialisierung & Die
Übergabe prüft vollständige Modelle und wählt den punktreichsten Kandidaten;
aktuelle Parameter wurden separat untersucht. & Modellwahl und
Registrierungsquote sind Teil des Datenvertrags.\\
Bilddomäne & Masken und Kameras beschrieben unterschiedliche Pixelräume & Das
Originalmodell bleibt erhalten; Bilder und Masken werden gemeinsam in eine
ideale PINHOLE-Domäne überführt. & Semantische und geometrische Daten müssen in
derselben Bilddomäne liegen.\\
Eval-Split & Leere Testkameraliste trotz 240 gültiger Maskenkameras & Namen wie
\texttt{00000} und \texttt{00000.jpg} wurden als normalisierte Stems und
numerische IDs abgeglichen; leere Listen werden abgelehnt. & Splitvalidierung
muss vor Training und Metrikberechnung erfolgen.\\
Matrixrunner & Nur die erste Variante wurde abgearbeitet & Unterprozesse
konsumierten den Standardinput der TSV-Schleife; der Runner verwendet nun
stdin-unabhängige Schleifen. & Ein Dry-Run ersetzt nicht die Kontrolle der
tatsächlich archivierten Läufe.\\
Metrik & Ein ausgewählter Eval-Frame besaß eine leere Maske & Der feste Split
wird nach dem Maskenwarp ausschließlich aus nichtleeren idealen Masken erzeugt.
& Globale Masken-Coverage und konkrete Eval-Eignung sind getrennte Kriterien.\\
SuGaR-Render & Importfehler und falsches Tensorlayout & Der Containerpfad wird
explizit importiert; HWC-Ausgaben werden vor dem Speichern in CHW überführt. &
Vorhandene STS-Metriken sind keine erfolgreichen SuGaR-Metriken.\\
Poisson & Leere oder ungültige Background-Punktwolke & Zu kleine Clouds und
ungültige Normalen werden übersprungen, solange die Foreground-Geometrie gültig
ist. & Diagnosekomponenten dürfen den Objektpfad nicht fälschlich abbrechen.\\
Centerline & Der Netzwerkmodus verwarf alle kurzen Äste & Der
\texttt{single}-Modus verwendet den Durchmesserpfad der größten Komponente. &
Verzweigte Netze benötigen eine getrennte wissenschaftliche Untersuchung.\\
\shortstack{Georeferenz-\\ierung} & Keine validierte 4x4-Matrix im Archivlauf & Der Export wird
als Translation-Fallback gekennzeichnet. & Ein technisch erzeugtes GeoJSON ist
kein Genauigkeitsnachweis.\\
\bottomrule
\end{tabularx}
\caption{Ausgewählte Fehlerfälle, Korrekturen und übertragbare Erkenntnisse.}
\label{tab:fehlerchronik}
\end{table}

Vollständige Laufprotokolle, Tracebacks und Zwischenartefakte liegen in den
Run-Archiven unter \texttt{data/10\_runs/}. Historische Sonnenbrillen- und
Gestellablationen bleiben Material zur Entwicklung der SuGaR-Route, werden aber
nicht als Ergebnisse des aktuellen Alurohrdatensatzes ausgegeben.

Die zu den Korrekturen gehörenden visuellen Zwischenstände sind über die
historischen Screenshot-Panels im Haupttext belegt: die STS- und
Objektfilterungsentwicklung in Kapitel~\ref{sec:implementierung}, die
SuGaR-Coarse-Zwischenstände in Abschnitt~\ref{sec:sugar-folgepruefung}. Die
Rohdateien liegen im Repository unter \texttt{PA/Screenshots/}; im Anlagenband
sind die zusammengeführten Panels unter \texttt{06\_Panels/} abgelegt (siehe
\texttt{ABGABE\_Index.md}).
